%! Absolute Value Equations & Inequalities — Unit 1, Lesson 1.2 — Guided Notes (Answer Key)
% Mirrors notes/main.tex exactly; every \blank{} becomes an \ans{} and every work block is
% authored byte-identically, so blank and key paginate alike.
\documentclass[10pt]{article}
\usepackage{algebra2-article}
\usepackage{algebra2-key}   % pulls in -boxes; do NOT also load -boxes
% Plain number line: {xmin}{xmax}, ticks labelled every unit.
\newcommand{\numline}[2]{%
  \draw[->, semithick, charcoal] (#1-0.5,0)--(#2+0.5,0);
  \foreach \x in {#1,...,#2}{\draw[charcoal] (\x,0.12)--(\x,-0.12) node[below, font=\scriptsize]{$\x$};}}
% The parkway: mile markers 0..20, ticks every mile, labels every four, tower at mile 12.
\newcommand{\trail}{%
  \draw[->, semithick, charcoal] (-0.7,0)--(21.2,0) node[right, font=\scriptsize]{mile};
  \foreach \x in {0,...,20}{\draw[charcoal] (\x,0.20)--(\x,-0.20);}
  \foreach \x in {0,4,...,20}{\node[below, font=\scriptsize, charcoal] at (\x,-0.24){$\x$};}
  \fill[forest] (12,0) circle (2.8pt);
  \node[forest, font=\scriptsize, above] at (12,0.28){tower};}
% Vocab answer for the key: bold forest term, then the filled definition on a write-line.
% The leading and trailing \par are required: \ansline ends with \dotfill but does NOT end the
% paragraph, so without them each term label gets pulled onto the previous answer's line.
\newcommand{\vocabans}[2]{%
  \par\noindent\textbf{\textcolor{forest}{#1:}}\\[1pt]\ansline{#2}\par}

\begin{document}

\pageheader{Unit 1, Lesson 1.2}{Guided Notes --- Answer Key}

\vspace{0.10in}

\begin{objectivebox}
\textbf{By the end of this lesson, I will be able to\dots}
\begin{itemize}\setlength{\itemsep}{2pt}
  \item read $|x-h|$ as the \emph{distance} between $x$ and $h$, and explain why an absolute value
        equation usually has \emph{two} solutions;
  \item \emph{isolate} the bars, then split into two cases, and handle the $c=0$ and $c<0$
        special cases;
  \item solve $|X|<c$ and $|X|>c$ and say which gives \emph{one} interval and which gives
        \emph{two} rays;
  \item write a solution set three ways --- set notation, interval notation, and a number line.
\end{itemize}
\end{objectivebox}

\vspace{0.05in}

\begin{vocabbox}
\small
Fill in each term as we build it together.
\par\vspace{2pt}   % \par is required: \vocabans's \noindent is a no-op mid-paragraph
\vocabans{Absolute value as distance}{$|x-h|$ is how far $x$ is from $h$; never negative}
\vocabans{Two-case rule}{$|X|=c$ with $c>0$ splits into $X=c$ or $X=-c$ --- isolate the bars first}
\vocabans{Compound ``and'' / ``or''}{$|X|<c$ gives one interval between; $|X|>c$ gives two rays outside}
\vocabans{Set notation}{$\{x \mid 7 \le x \le 17\}$ --- ``all $x$ such that\dots''}
\vocabans{Interval notation}{$[7,17]$; brackets include an endpoint, parentheses exclude it}
\vocabans{Extraneous / no solution}{$|X|=c$ or $|X|<c$ with $c<0$ has no solution --- distance is never negative}
\end{vocabbox}

\begin{hookbox}
A straight stretch of parkway is marked off in miles, $0$ through $20$. A ranger's radio tower
stands at \textbf{mile 12}, and a handheld radio reaches it only from \textbf{within 5 miles}.
\begin{center}
\begin{tikzpicture}[x=0.52cm, y=0.5cm]
  \trail
\end{tikzpicture}
\end{center}
\vspace{-2pt}
\begin{itemize}\setlength{\itemsep}{1pt}
  \item Two hikers stand \emph{exactly} $5$ miles from the tower. Their mile markers are
        \ans{$7$} and \ans{$17$}.
  \item Why are there \emph{two} of them, and not one? \ans{one on each side of the tower}
\end{itemize}
Every question today is this picture: how far is $x$ from the tower --- and is that close enough?
\end{hookbox}

\boxguard[22]
\begin{notesbox}{1. Absolute Value Is a Distance --- So It Answers Twice}
\begin{center}
\begin{minipage}[c]{0.50\linewidth}
\[ |x-h| \;=\; \text{the distance between } x \text{ and } h \]
On the parkway, a hiker at mile $x$ is \ans{$|x-12|$} miles from the tower. A distance is never
\ans{negative}, and \emph{two} different spots sit the same distance out --- one on each side.
\end{minipage}\hfill
\begin{minipage}[c]{0.46\linewidth}\centering
\begin{tikzpicture}[x=0.30cm, y=0.42cm]
  \trail
  \draw[very thick, redacc] (7,0.55)--(12,0.55);
  \draw[very thick, redacc] (12,0.55)--(17,0.55);
  \fill[redacc] (7,0) circle (2.8pt);
  \fill[redacc] (17,0) circle (2.8pt);
  \node[redacc, font=\tiny, above] at (9.5,0.6){$5$};
  \node[redacc, font=\tiny, above] at (14.5,0.6){$5$};
\end{tikzpicture}
\end{minipage}
\end{center}
\textbf{Worked: exactly five miles out.} ``Distance from the tower is $5$'' is $|x-12|=5$.
\begin{work}
  |x - 12| &= 5 \\
  x - 12 &= 5 \quad\text{or}\quad x - 12 = -5 \\
  x &= 17 \quad\text{or}\quad x = 7
\end{work}
Both markers from the hook, straight out of the algebra. \textbf{Verify} by putting each back in:
$|17-12|$ is \ans{$5$} and $|7-12|$ is \ans{$5$}.
\end{notesbox}

\boxguard[24]
\begin{notesbox}{2. Isolate the Bars \emph{Before} You Split}
The two-case rule applies to $|X| = c$ only when the bars stand \emph{alone}. Everything outside
them gets undone first --- exactly as if the bars were a variable.
\begin{center}
\renewcommand{\arraystretch}{1.5}
\begin{tabularx}{\linewidth}{@{}l l X@{}}
\textbf{$|X| = c$ when\dots} & \textbf{Split into} & \textbf{How many solutions --- and why} \\
\hline
$c > 0$ & $X =$ \ans{$c$} \; or \; $X =$ \ans{$-c$} & \ans{two} --- two spots sit that far out \\
$c = 0$ & $X =$ \ans{$0$} & \ans{one} --- only one spot is \emph{zero} away \\
$c < 0$ & --- & \ans{none} --- a distance is never negative \\
\end{tabularx}
\end{center}
\textbf{Worked: undo, then split.} Solve $3|x+1| - 2 = 10$.
\begin{work}
  3|x+1| - 2 &= 10 \\
  3|x+1| &= 12 \\
  |x+1| &= 4 \\
  x + 1 &= 4 \quad\text{or}\quad x + 1 = -4 \\
  x &= 3 \quad\text{or}\quad x = -5
\end{work}
Splitting at the \emph{first} line would give $3x+1-2=10$ --- a single answer, and a wrong one.
\end{notesbox}

\boxguard[26]
\begin{notesbox}{3. ``Less th\emph{AND}'' vs.\ ``Great\emph{OR}'' --- One Piece or Two}
Same expression, same $5$, one symbol apart --- and the two answers look nothing alike.
\begin{center}
\renewcommand{\arraystretch}{1.4}
\begin{tabularx}{\linewidth}{@{}X|X@{}}
\textbf{In range: $|x-12| \le 5$} & \textbf{Out of range: $|x-12| \ge 5$} \\
\hline
Closer to the tower than $5$: \emph{between}. & Farther from the tower than $5$: \emph{outside}. \\
Becomes one compound \ans{and}: & Becomes two statements joined by \ans{or}: \\
$-5 \le x-12 \le 5$, so \; \ans{$7$} $\le x \le$ \ans{$17$} & $x-12 \le -5$ \; or \; $x-12 \ge 5$, so \; $x \le$ \ans{$7$} \, or \, $x \ge$ \ans{$17$} \\
\centering\arraybackslash
\begin{tikzpicture}[x=0.20cm, y=0.40cm, baseline=-2pt]
  \trail
  \draw[very thick, greenacc] (7,0)--(17,0);
  \fill[greenacc] (7,0) circle (3.6pt);
  \fill[greenacc] (17,0) circle (3.6pt);
\end{tikzpicture}
&
\centering\arraybackslash
\begin{tikzpicture}[x=0.20cm, y=0.40cm, baseline=-2pt]
  \trail
  \draw[very thick, redacc] (0,0)--(7,0);
  \draw[very thick, redacc] (17,0)--(20,0);
  \fill[redacc] (7,0) circle (3.6pt);
  \fill[redacc] (17,0) circle (3.6pt);
\end{tikzpicture}
\\
\ans{one} connected piece & \ans{two} separate pieces \\
\end{tabularx}
\end{center}
\vspace{2pt}
\begin{tcolorbox}[colback=goldbg, colframe=goldacc, boxrule=0.9pt, arc=2mm,
    left=2.5mm, right=2.5mm, top=1.5mm, bottom=1.5mm]
\small \textbf{Careful --- test one hiker and watch the two disagree.} A hiker at \textbf{mile 3}
is $|3-12| =$ \ans{$9$} miles from the tower. Is $9 \le 5$? \ans{no} \; Is $9 \ge 5$?
\ans{yes} \; So mile $3$ belongs to \emph{exactly one} of the two answers above ---
\ans{out of range}. Picking the wrong symbol does not shift the answer a little; it hands you the
\emph{opposite} set.
\end{tcolorbox}
\end{notesbox}

\boxguard[24]
\begin{notesbox}{4. Three Ways to Write the Same Set}
The in-range hikers are $7 \le x \le 17$. That one set has three accepted spellings.
\begin{center}
\renewcommand{\arraystretch}{1.5}
\begin{tabularx}{\linewidth}{@{}l X@{}}
\textbf{Set notation} & $\{x \mid$ \ans{$7 \le x \le 17$} $\}$ --- read ``all $x$ such that\dots'' \\
\textbf{Interval notation} & \ans{$[7,17]$} --- endpoints in order, left to right \\
\textbf{Number line} &
\begin{tikzpicture}[x=0.20cm, y=0.40cm, baseline=-2pt]
  \trail
  \draw[very thick, keyred] (7,0)--(17,0);
  \fill[keyred] (7,0) circle (3.6pt);
  \fill[keyred] (17,0) circle (3.6pt);
\end{tikzpicture} \\
\end{tabularx}
\end{center}
\vspace{-2pt}
\begin{center}
\renewcommand{\arraystretch}{1.4}
\begin{tabularx}{\linewidth}{@{}X|X@{}}
\textbf{Endpoint included} ($\le$, $\ge$) & \textbf{Endpoint excluded} ($<$, $>$) \\
\hline
bracket \ans{$[\;]$} \, and a \ans{filled} dot & parenthesis \ans{$(\;)$} \, and an \ans{open} dot \\
\end{tabularx}
\end{center}
Two rays need two intervals joined by $\cup$: the out-of-range hikers are
$(-\infty,$ \ans{$7$} $] \cup [$ \ans{$17$} $, \infty)$. Next to $\infty$ you always write a
\ans{parenthesis}, never a bracket --- because $\infty$ is not a \ans{number} you can reach.
\end{notesbox}

\begin{practicebox}
Solve $2|x-5| - 3 = 7$, then change one symbol and watch the answer change shape.
\begin{enumerate}[leftmargin=1.7em, itemsep=3pt]
  \item Isolate the bars, then split.
        \begin{work}
          2|x-5| - 3 &= 7 \\
          2|x-5| &= 10 \\
          |x-5| &= 5 \\
          x - 5 &= 5 \quad\text{or}\quad x - 5 = -5 \\
          x &= 10 \quad\text{or}\quad x = 0
        \end{work}
  \item Now solve $2|x-5| - 3 \le 7$ instead. One piece or two? \ans{one} \;
        Solution: \ans{$0 \le x \le 10$}
  \item And $2|x-5| - 3 \ge 7$? \ans{two} \; Solution: \ans{$x \le 0$ or $x \ge 10$}
\end{enumerate}
\end{practicebox}

\boxguard[22]
\begin{scenariobox}[Individual Practice --- On Your Own]{navy}
Three problems, quietly and on your own. Isolate the bars \emph{before} you split, and decide each
time whether the answer is one piece or two.
\begin{enumerate}[leftmargin=1.9em, itemsep=9pt]
  \item \textbf{An equation.} Solve $4|x+1| - 5 = 15$, then verify one of your two answers.
    \begin{work}
      4|x+1| - 5 &= 15 \\
      4|x+1| &= 20 \\
      |x+1| &= 5 \\
      x + 1 &= 5 \quad\text{or}\quad x + 1 = -5 \\
      x &= 4 \quad\text{or}\quad x = -6
    \end{work}
    Solutions: \ans{$4$ and $-6$} \quad Check one: \ans{$x=4$: $4(5)-5=15$ \checkmark}

  \item \textbf{``Less th\emph{AND}.''} Solve $|2x - 1| \le 5$, then write the set three ways.
    \begin{work}
      |2x - 1| &\le 5 \\
      -5 \le 2x - 1 &\le 5 \\
      -4 \le 2x &\le 6 \\
      -2 \le x &\le 3
    \end{work}
    Set: \ans{$\{x \mid -2 \le x \le 3\}$} \quad Interval: \ans{$[-2,3]$} \quad
    \begin{tikzpicture}[x=0.40cm, y=0.42cm, baseline=-2pt]
      \numline{-4}{5}
      \draw[very thick, keyred] (-2,0)--(3,0);
      \fill[keyred] (-2,0) circle (2.8pt);
      \fill[keyred] (3,0) circle (2.8pt);
    \end{tikzpicture}

  \item \textbf{``Great\emph{OR}.''} Solve $|3x - 6| \ge 9$, then write the set three ways.
    \begin{work}
      |3x - 6| &\ge 9 \\
      3x - 6 &\le -9 \quad\text{or}\quad 3x - 6 \ge 9 \\
      3x &\le -3 \quad\text{or}\quad 3x \ge 15 \\
      x &\le -1 \quad\text{or}\quad x \ge 5
    \end{work}
    Set: \ans{$\{x \mid x \le -1$ or $x \ge 5\}$} \quad Interval: \ans{$(-\infty,-1] \cup [5,\infty)$} \quad
    \begin{tikzpicture}[x=0.40cm, y=0.42cm, baseline=-2pt]
      \numline{-3}{7}
      \draw[very thick, keyred] (-3.5,0)--(-1,0);
      \draw[very thick, keyred] (5,0)--(7.5,0);
      \fill[keyred] (-1,0) circle (2.8pt);
      \fill[keyred] (5,0) circle (2.8pt);
    \end{tikzpicture}
\end{enumerate}
\end{scenariobox}

\end{document}
